%% kapa: exam=1 round=34 year=2023 period=2 form=A subject=appraisal_law subject_ko=감정평가관계법규 seq=35 number=35
%% source: Q-Net 감정평가사 기출문제 (cst00309) artlSeq=5211789 file=raw/1008/5211789/2023년도 제34회 감정평가사 1차시험 2교시.pdf
%% conversion: kordoc 4.24.51 -> tools/gen1.py (auto); review=proofread-sonnet; flags=none
\begin{question}{35}
부동산등기법령상 소유권등기 및 담보권등기에 관한 설명으로 옳지 않은 것은?
\begin{choices}
  \item 등기관이 소유권보존등기를 할 때에는 등기원인과 그 연월일을 기록하지 아니한다.
  \item 미등기의 토지 또는 건물의 경우 수용(收用)으로 인하여 소유권을 취득하였음을 증명하는 자는 소유권보존등기를 신청할 수 있다.
  \item 등기관이 동일한 채권에 관하여 3개의 부동산에 관한 권리를 목적으로 하는 저당권설정의 등기를 할 때에는 공동담보목록을 작성하여야 한다.
  \item 등기관은 근저당권을 내용으로 하는 저당권설정의 등기를 할 때에는 채권의 최고액을 기록하여야 한다.
  \item 등기관이 채권의 일부에 대한 대위변제로 인한 저당권 일부이전등기를 할 때에는 변제액을 기록하여야 한다.
\end{choices}
\end{question}
\answer{3}
